%! TeX root = Definición_Planeteamiento_Proyecto.tex

\begin{frame}{Experimental Framework}
\blurbackground % Activates your custom blur effect
\footnotesize
\centering
\begin{adjustbox}{max width=\textwidth}
\begin{tabular}{p{1.2cm}p{2.2cm}p{2.2cm}p{1.8cm}p{2cm}p{2cm}}
\toprule
\textbf{Exp ID} & \textbf{Research Question} & \textbf{Methodology} & \textbf{Input Parameters} & \textbf{Success Metrics} & \textbf{Expected Outcome} \\
\midrule
\textbf{EXP-1} & Can we reproduce the 2020 paper's null result? & Implement exact 2020 model. Use GA + ANN. & 14 flux integers, 6 moduli fields & \% critical points, Stability ratio & Confirm absence of stable dS vacua \\
\midrule
\textbf{EXP-2} & Does torsion geometry enable stable dS states? & Implement 2022 model. Use Hybrid SA+CG. & 5 flux coefficients, 2 moduli & Number of stable dS, Flatness measure & Verify claim of ~170 stable dS vacua \\
\midrule
\textbf{EXP-3} & Which ML strategy finds viable vacua fastest? & Compare: GA, SA, NN, PSO on both models & Identical parameter spaces & Convergence speed, Solution quality & Identify optimal algorithm \\
\midrule
\textbf{EXP-4} & How sensitive are "stable" dS vacua? & Apply perturbations: quantum corrections, flux variations & Perturbation strength parameters & Lifetime estimation, Decay channels & Classify stability vs metastability \\
\midrule
\textbf{EXP-5} & Where are Swampland boundaries? & Systematic scan near conjectured boundaries & c' parameter, moduli mass ratios & Distribution in (m²/V) space & Precisely map Swampland boundary \\
\midrule
\textbf{EXP-6} & How do results scale with complexity? & Gradually increase: moduli, fluxes, geometries & Dimensionality parameters & Computational scaling laws & Understand realistic model limitations \\
\bottomrule
\end{tabular}
\end{adjustbox}
\end{frame}

% Additional supporting frames for your thesis proposal:

\begin{frame}{Research Hypothesis \& Novel Contributions}
\blurbackground
\begin{tcolorbox}[
    enhanced,
    colback=black!85,
    colframe=Magenta,  % Using your magenta color for emphasis
    arc=3mm,
    boxrule=1.5pt,
    title={\color{Yellow!80!white}Hypothesis}
]
\color{White}
The claimed stable de Sitter vacua in torsional compactifications are actually \textbf{metastable} when subjected to robust stability analysis, and advanced ML algorithms can precisely quantify their lifetime and decay channels.
\end{tcolorbox}

\vspace{0.5cm}

\begin{tcolorbox}[
    enhanced,
    colback=black!85,
    colframe=LimeGreen,  % Using your lime green color
    arc=3mm,
    boxrule=1.5pt,
    title={\color{Yellow!80!white}Novel Contributions}
]
\begin{itemize}
    \item First systematic comparison of ML methods for string landscape exploration
    \item Quantitative stability analysis of claimed dS vacua
    \item Empirical mapping of Swampland conjecture boundaries
    \item Scalability analysis for computational string theory
\end{itemize}
\end{tcolorbox}
\end{frame}

\begin{frame}{Methodology Overview}
\blurbackground
\begin{columns}
\begin{column}{0.48\textwidth}
\begin{tcolorbox}[
    enhanced,
    colback=black!85,
    colframe=mypurple,
    arc=3mm,
    boxrule=1pt,
    title={\color{Yellow!80!white}Algorithms to Compare}
]
\begin{itemize}
    \item Genetic Algorithm (GA)
    \item Simulated Annealing (SA)
    \item Neural Networks (NN)
    \item Particle Swarm (PSO)
    \item Hybrid SA+CG
\end{itemize}
\end{tcolorbox}
\end{column}

\begin{column}{0.48\textwidth}
\begin{tcolorbox}[
    enhanced,
    colback=black!85,
    colframe=mypurple,
    arc=3mm,
    boxrule=1pt,
    title={\color{Yellow!80!white}Evaluation Metrics}
]
\begin{itemize}
    \item Convergence speed
    \item Solution quality (V\textsubscript{min})
    \item Computational cost
    \item Stability classification
    \item Boundary precision
\end{itemize}
\end{tcolorbox}
\end{column}
\end{columns}

\vspace{0.5cm}

\begin{center}
\begin{tcolorbox}[
    enhanced,
    colback=black!85,
    colframe=Yellow!80!white,
    arc=3mm,
    boxrule=1pt,
    width=0.8\textwidth,
    title={\color{mypurple}Technical Implementation}
]
\begin{itemize}
    \item \texttt{Python} (TensorFlow/PyTorch for ML)
    \item High-performance computing cluster
    \item Parallel evaluation of flux configurations
    \item Automated stability analysis pipeline
\end{itemize}
\end{tcolorbox}
\end{center}
\end{frame}

\begin{frame}{Deliverables}
\blurbackground
\begin{tcolorbox}[
    enhanced,
    colback=black!85,
    colframe=AntiqueWhite,
    arc=3mm,
    boxrule=1pt,
    title={\color{mypurple}Expected Deliverables}
    ]
    \begin{itemize}
        \item Comparative analysis framework
        \item Stability classification system
        \item Boundary mapping methodology
        \item Open-source code repository
        \item Thesis document
    \end{itemize}
\end{tcolorbox}

\begin{tcolorbox}[
    enhanced,
    colback=black!85,
    colframe=Red!80!white,
    arc=3mm,
    boxrule=1pt,
    title={\color{Yellow!80!white}Success Criteria}
    ]
    \begin{itemize}
        \item Reproduce published results
        \item Identify optimal algorithms
        \item Quantify stability bounds
        \item Provide scalable methodology
    \end{itemize}
\end{tcolorbox}
\end{frame}
